
\item \textbf{Выбор уставки токового органа блокировки РПН}

%\begin{enumerate}[label=\arabic{section}.\arabic{enumi}.\arabic*, labelsep=4pt, leftmargin=0em, itemindent=65pt, parsep=3pt]
\begin{enumerate}[label={}, labelsep=4pt, leftmargin=0pt, itemindent=57pt]
\item
Т.к. тип устройства РПН не известен, то принимаем его номинальный ток равным номинальному току стороны, на которой он установлен (то есть сторона ВН).

По выражению (\ref{eq:torpn2}) определяется первичный ток срабатывания блокирующего органа РПН $I_{\textsl{бл.РПН}}$:
\begin{equation*}
    I_{\textsl{бл.РПН}} = {K_{\textsl{отс}}}\cdot {K_{\textsl{пер}}}\cdot {I_{\textsl{ном.ст.РПН}}} = 0,95 \cdot 2,0 \cdot 250 = 475\; (\textsl{А}),
\end{equation*}
Уставка задается в относительных единицах от номинального тока аналогового входа устройства, для этого полученное значение пересчитывается по следующему выражению:
\begin{equation*}
    I_{\textsl{ср}} = \frac{I_{\textsl{бл.РПН}}}{I_{\textsl{ном.вх.уст}}} \cdot \frac{I_{\textsl{ном.вт.ТТ}}}{I_{\textsl{ном.пер.ТТ}}} = \frac{475}{5} \cdot \frac{5}{400} = 1,2\; (\textsl{о.е.}),
\end{equation*}
\begin{eqexpl}[15mm]
    \item{$I_{\textsl{ном.вх.уст}}$} номинальный ток аналоговых входов устройства, 1 или 5 А;
    \item{$I_{\textsl{ном.вт.ТТ}}$} номинальный вторичный ток ТТ, А;
    \item{$I_{\textsl{ном.пер.ТТ}}$} номинальный первичный ток ТТ, А.
\end{eqexpl}

Параметр \textbf{$I_{\textsl{ср}}$} <<Ток срабатывания>> функции ТО РПН принимается \textbf{равным 1,2}.
\end{enumerate}
